\subsection{与模相关的循环}

本节中，A 表示一个诺特环。

{}{}设 Z(A) 是以 Spec(A) 的不可约闭子集集合为基的自由 Z-模；对于 Spec(A) 的每个不可约闭子集 Y，记 [Y] 为 Z(A) 中相应的元素。Z(A) 的元素有时称为循环。

设 M 是一个有限生成 A-模。对于 A 的每个素理想 只要它是  的极小元，就有 \(0 < \text{long}_{A_{p}}(M_{p}) < \infty\)（IV, § 2, no. 5, 命题 7 的推论 2 以及 § 1, no. 4, theorem 2）；令

\[
z (\mathbf {M}) = \sum_ {\mathfrak {p}} \operatorname{long} _ {\mathrm{A} _ {\mathfrak {p}}} (\mathrm{M} _ {\mathfrak {p}}). [ \mathrm{V} (\mathfrak {p}) ],
\]

其中  遍历 Supp(M) 的极小素理想组成的有限集合。

注。--- 对每个 \(\mathfrak{p} \in \operatorname{Spec}(\mathbf{A})\)，有 \(z(\mathbf{A}/\mathfrak{p}) = [\mathbf{V}(\mathfrak{p})]\)。更一般地，设 M 是一个有限生成 A-模，且 \((\mathbf{M}_i)_{0 \leqslant i \leqslant n}\) 是 M 的一个合成列，使得对 \(0 \leqslant i \leqslant n - 1\)，模 \(\mathbf{M}_i / \mathbf{M}_{i+1}\) 同构于 \(\mathbf{A}/\mathfrak{p}_i\)，其中 \(\mathfrak{p}_i\) 是 A 的一个素理想（参见 IV, § 1, no. 4, theorem 1）；则有 \(z(\mathbf{M}) = \sum_{i \in \mathbf{J}} [\mathbf{V}(\mathfrak{p}_i)]\)，其中  是由满足 \(i\) 且 \(\mathfrak{p}_i\) 是 \(\{\mathfrak{p}_0, \ldots, \mathfrak{p}_{n-1}\}\) 的极小元的那些 i 组成的 I 的子集（IV, § 1, no. 4, theorem 2 以及 § 2, no. 5, 注 1）。

对每个整数 \(d\)，记 \(Z_{\leq d}\)（分别为 \(Z_d\)，分别为 \(Z^{\geq d}\)，分别为 \(Z^d\)）为子-\(\mathbf{Z}\)-模 \(Z(A)\) 中由满足下述条件的元素  生成的模：其中 \([V(\mathfrak{p})]\) 是 A 的一个素理想，且 \(\mathfrak{p}\)（分别为 \(A\)，分别为 \(\dim(A/\mathfrak{p}) \leqslant d\)，分别为 \(\dim(A/\mathfrak{p}) = d\)\(\operatorname{ht}(\mathfrak{p}) \geqslant d\)\(\operatorname{ht}(\mathfrak{p}) = d\)）。\(Z_{d}\)（分别为 \(Z^{d}\)）的元素称为维数为 \(d\) 的循环（分别称为余维数为 \(d\) 的循环）。显然有

\[
\mathbf {Z} _ {\leqslant d} = \mathbf {Z} _ {\leqslant d - 1} \oplus \mathbf {Z} _ {d}, \quad \mathbf {Z} ^ {\geqslant d} = \mathbf {Z} ^ {\geqslant d + 1} \oplus \mathbf {Z} ^ {d}.
\]

令 C 为有限生成 A-模的类的集合（A, VIII, § 3, no. 5），对每个整数 d，令 \(C_{\leq d}\)（分别为 \(C^{\geq d}\)）为 C 的子集，由维数 \(\leq d\) 的有限生成 A-模的类（分别为其支集在 Spec(A) 中余维数 \(\geq d\) 的有限生成 A-模的类）组成。

引理 1. --- 设 M 是一个有限生成 A-模，d 是一个整数。

\begin{enumerate}
\def\labelenumi{\alph{enumi})}
\tightlist
\item
  要使 M 属于 \(C_{\leq d}\)，必要且充分条件是 \(z(\mathbf{M}) \in \mathbf{Z}_{\leq d}\)；此时，\(z_{d}(\mathbf{M})\) 是将 \(z(\mathbf{M})\) 投影到 \(Z_{d}\)（沿与 \(Z_{\leq d-1}\) 平行的方向）所得的投影，下面给出其表达式：
\end{enumerate}

\[
z _ {d} (\mathrm{M}) = \sum _ {\dim (\mathrm{A} / \mathfrak {p}) = d} \operatorname{long} _ {\mathrm{A} _ {\mathfrak {p}}} (\mathrm{M} _ {\mathfrak {p}}). [ \mathrm{V} (\mathfrak {p}) ].
\]

\begin{enumerate}
\def\labelenumi{\alph{enumi})}
\setcounter{enumi}{1}
\tightlist
\item
  要使 M 属于 \(C^{\geq d}\)，必要且充分条件是 \(z(\mathbf{M})\in\mathbf{Z}^{\geq d}\)；此时，\(z^{d}(\mathbf{M})\) 是将 \(z(\mathbf{M})\) 投影到 \(\mathbf{Z}^{d}\)（沿与 \(\mathbf{Z}^{\geq d+1}\) 平行的方向）所得的投影，下面给出其表达式：
\end{enumerate}

\[
z ^ {d} (\mathbf {M}) = \sum _ {\operatorname{ht} (\mathfrak {p}) = d} \operatorname{long} _ {\mathrm{A} _ {\mathfrak {p}}} (\mathrm{M} _ {\mathfrak {p}}). [ \mathrm{V} (\mathfrak {p}) ].
\]

要使 M 属于 \(C_{\leq d}\)，即维数 \(\leq d\)，必要且充分条件是：对于  的每个极小素理想  ，有 \(\leq d\)，这意味着 \(z(M) \in Z_{\leq d}\)。假设  ，并令 \(\leq d\) 且 \(p \in Spec(A)\)；则或者 \(p \notin Supp(M)\)，因而 \(M_p = 0\)，或者 \(p \in Supp(M)\)，且 { 是  的极小元 }；在这两种情形下，{} 中 \(z(M)\) 的系数都是 \(_{A_p}(M_p)\)，从而得到 a)。b) 类似可证；注意，有限生成模 M 属于 \(C^{\geq d}\)，当且仅当对每个素理想 \(M_{\mathfrak{p}} = 0\) 的高度 \textless{}。

根据 no. 4 的命题 9 c) 和注 3，集合 \(\mathcal{C}_{\leq d}\) 及 \(\mathcal{C}^{\geq d}\) 是遗传的（A, VIII, § 10, no. 1, 定义 1），因而可以考虑与它们对应的 Grothendieck 群  和  （同上，no. 2）；对于 \(\mathcal{C}_{\leq d}\)（分别为 \(\mathcal{C}^{\geq d}\)）中的每个有限型 A-模 M，我们将  （分别为 \(\mathcal{C}_{\leq d}\)\(\mathcal{C}^{\geq d}\){）的关联元素记为 }{}\(_{\leq d}\)（分别为 {}{}\(^{\geq d}\)\(\mathcal{C}_{\leq d}\)\(\mathcal{C}^{\geq d}\)）。由引理 1，函数 \(z_{d}\) 和 \(z^{d}\) 是可加的；因此存在（同上，命题 3）同态

\[
\zeta_ {d}: \mathrm{K} (\mathbb {C} _ {\leqslant d}) \rightarrow Z _ {d}, \quad \zeta^ {d}: \mathrm{K} (\mathbb {C} ^ {\geqslant d}) \rightarrow Z ^ {d}
\]

使得 \(\zeta_d([M]_{\leq d}) = z_d(M)\) 对每个 \(C_{\leq d}\) 型的有限生成 A-模成立，并且 \(\zeta^d ([N]^{\geq d}) = z^d (N)\) 对每个 \(C^{\geq d}\) 型的有限生成 A-模 N 成立。此外，由于 \(C_{\leq d-1} \subset C_{\leq d}\) 且 \(C^{\geq d+1} \subset C^{\geq d}\)，我们有典范同态

\[
i _ {d}: \mathrm{K} (\mathbb {C} _ {\leqslant d - 1}) \rightarrow \mathrm{K} (\mathbb {C} _ {\leqslant d}) \quad \text {\text{ 与 }} \quad i ^ {d}: \mathrm{K} (\mathbb {C} ^ {\geqslant d + 1}) \rightarrow \mathrm{K} (\mathbb {C} ^ {\geqslant d}).
\]

用这些记号：

命题 10. --- Z-模及同态序列

\[
\mathrm{K} \left(\mathbb {C} _ {\leqslant d - 1}\right) \xrightarrow {i _ {d}} \mathrm{K} \left(\mathbb {C} _ {\leqslant d}\right) \xrightarrow {\zeta_ {d}} \mathrm{Z} _ {d} \longrightarrow 0
\]

\[
\mathrm{K} \left(\mathbb {C} ^ {\geqslant d + 1}\right) \xrightarrow {i ^ {d}} \mathrm{K} \left(\mathbb {C} ^ {\geqslant d}\right) \xrightarrow {\zeta ^ {d}} \mathrm{Z} ^ {d} \longrightarrow 0
\]

是正合的。

我们有 \(\zeta_d \circ i_d = 0\)。根据引理 1，对于每个满足 \(\mathfrak{p} \in \operatorname{Spec}(A)\) 且 \(\dim(A/\mathfrak{p}) = d\) 的素理想，有 \(\zeta_d([A/\mathfrak{p}]_{\leq d}) = z_d(A/\mathfrak{p}) = [V(\mathfrak{p})]\)，因而同态 \(\zeta_d\) 是满射。根据 IV, § 1, no. 4, th. 1，\(K(C_{\leq d})\) 由 \([A/\mathfrak{p}]_{\leq d}\) 生成，其中 \(\mathfrak{p} \in \operatorname{Spec}(A)\) 且 \(\dim(A/\mathfrak{p}) \leqslant d\)；因此，每个元素 \(\xi\) 属于 \(K(C_{\leq d})\) 都可写成 \(\xi = i_d(\eta) + \sum_{i=1}^{k} n_i[A/\mathfrak{p}_i]_{\leq d}\)，其中 \(\eta \in K(C_{\leq d-1})\)、\(n_i \in Z\)，且 \(\dim(A/\mathfrak{p}_i) = d\) 对 \(1 \leqslant i \leqslant k\) 成立；有 \(\zeta_d(\xi) = \sum_{i=1}^{k} n_i[V(\mathfrak{p}_i)]\)，因而 \(\zeta_d(\xi) = 0\) 蕴含 \(\xi = i_d(\eta) \in \operatorname{Im}(i_d)\)，从而 \(\operatorname{Ker}(\zeta_d) = \operatorname{Im}(i_d)\)。

第二个序列同理可证。

例。--- 1) 假设 A 是诺特整环。则有 \(Z^0 = \mathbf{Z}\){}{}；有 \(C^{\geqslant 0} = C\) 且 \(z^0(M) = \operatorname{rg}(M)\)。[Spec(A)]。因此，\(C^{\geqslant 1}\) 型的模就是挠模。

\begin{enumerate}
\def\labelenumi{\arabic{enumi})}
\setcounter{enumi}{1}
\item
  假设 A 是诺特整闭环。则 \(Z^1\) 可识别为第七章（§ 1, no. 3, th. 2 以及 no. 6, th. 3）中引入的 A 的因子群 D(A)。\(C^{\geqslant 2}\) 型的模是伪零模（VII, § 4, no. 4, 定义 2）；如果 M 是一个有限生成挠模，则 \(z^1(M) \in Z^1 = D(A)\) 是 M 的内容 \(\chi(M)\)（VII, § 4, no. 5, 定义 4）。因此，同上所述的命题 10 和 11 等价于序列 \(K(C^{\geqslant 2}) \to K(C^{\geqslant 1}) \to Z^1 \to 0\) 的正合性。
\item
  \(C_{\leqslant0}\) 型的模就是维数 \(\leqslant0\) 的模，即有限长度模（no. 4, 注 1）。对于每个有限长度的 A-模 M，有 \(\text{long}_{\mathbf{A}}(\mathbf{M}) = \varepsilon(z_0(\mathbf{M}))\)，其中 \(\varepsilon: Z_0 \to Z\) 把线性组合 \(\sum_{m} n_{m}[V(m)]\) 赋予整数 \(\sum_{m} n_{m}\)（IV, § 2, no. 5, 命题 8 的推论）。
\item
  假设 A 是有限维整环。置 \(d = \dim(A)\)。则有 \(\mathcal{C}_{\leq d} = \mathcal{C}\)、\(Z_d = \mathbf{Z}\){}{}\(Z^0\)、\(z_d(M) = \operatorname{rg}(M)\)。[Spec(A)] = \(z^0(M)\)，而 \(\mathcal{C}_{\leq d-1}\) 型的模是挠模。
\end{enumerate}

